% =====================================================================================================
% theory.tex -- theory section for "planning with stochastic latent world models".
% Self-contained section (no \documentclass).  The host document must load amsmath, amssymb, amsthm
% (and a bibliography that knows the keys in theory_refs.bib).  Theorem environments are defined here
% only if the host has not defined them already.
% Numerical verification of every statement: experiments/d4_verify.py  (output: results/d4/).
% =====================================================================================================
\ifcsname theorem\endcsname\else\newtheorem{theorem}{Theorem}\fi
\ifcsname proposition\endcsname\else\newtheorem{proposition}[theorem]{Proposition}\fi
\ifcsname lemma\endcsname\else\newtheorem{lemma}[theorem]{Lemma}\fi
\ifcsname corollary\endcsname\else\newtheorem{corollary}[theorem]{Corollary}\fi
\ifcsname definition\endcsname\else\newtheorem{definition}[theorem]{Definition}\fi
\ifcsname remark\endcsname\else\newtheorem{remark}[theorem]{Remark}\fi
\ifcsname example\endcsname\else\newtheorem{example}[theorem]{Example}\fi
\providecommand{\EE}{\mathbb{E}}
\providecommand{\PP}{\mathbb{P}}
\providecommand{\RR}{\mathbb{R}}
\providecommand{\Var}{\operatorname{Var}}
\providecommand{\tr}{\operatorname{tr}}
\providecommand{\CVaR}{\operatorname{CVaR}}
\providecommand{\VaR}{\operatorname{VaR}}
\providecommand{\Bin}{\operatorname{Bin}}
\providecommand{\osc}{\operatorname{osc}}
\providecommand{\ES}{\operatorname{ES}}
\providecommand{\indic}{\mathbf 1}
\providecommand{\Nrm}{\mathcal{N}}
\providecommand{\Phibar}{\bar\Phi}

\section{Theory of planning with stochastic latent world models}\label{sec:theory}

This section collects the formal statements behind the empirical findings.
Every result is labelled \emph{Known} (standard, cited, proof recalled only when it is short) or \emph{New} (elementary, proved here;
``new'' means new in this form for the planning setting, not deep).
All formulas were checked numerically with \texttt{experiments/d4\_verify.py} (Section~\ref{sec:verification}).
Where the numerics forced a weaker statement than first conjectured, this is said explicitly (Remark~\ref{rem:firstpassage}).

\begin{table}[h]
\centering\footnotesize
\begin{tabular}{@{}lp{2.6cm}p{6.6cm}l@{}}
\hline
Result & Where & One-line statement & Status\\
\hline
T1 & Thm.~\ref{thm:regret} & regret of deterministic (mean-embedding) planning $=$ difference of noise premiums, $\le$ their spread, tight & New\\
   & Cor.~\ref{cor:fail}, \ref{cor:cvar} & failure-penalised cost (regret linear in $\kappa$); $\mathrm{CVaR}_\alpha$ (planner blind to $\alpha$) & New\\
T2 & Thm.~\ref{thm:select} & $\PP(\text{risky selected})=1-(1-F_{\mathrm{Bin}}(k^*;M,p))^{N_r}$, monotone in $N_r$, $\to1$ & New\\
   & Thm.~\ref{thm:excess} & excess true cost $\kappa(p-\theta)\PP(\text{risky})$; optimism $\to\kappa p$ & New\\
   & Thm.~\ref{thm:verify} & fresh-sample verification: risk $1-(1-F')^K$, free of $N_r$ & New\\
   & Prop.~\ref{prop:shrink} & monotone shrinkage with equal $M$ cannot change the argmin; unequal $M$ can & New\\
T3 & Thm.~\ref{thm:propriety} & energy score strictly proper & Known\\
   & Prop.~\ref{prop:noisemap} & law identified; noise map and cross-candidate coupling are not & New\\
   & Thm.~\ref{thm:consistency} & set-consistency of the empirical energy-score minimiser & New (standard proof)\\
T4 & Prop.~\ref{prop:var}, Thm.~\ref{thm:lqg} & variance $t\sigma^2$ vs bounded; open-loop value $=H\times$ closed-loop value & New\\
   & Thm.~\ref{thm:barrier} & strict open-loop gap with an absorbing boundary & New\\
T5 & Thm.~\ref{thm:scale} & no scale recalibration repairs a change of shape; explicit floors & New\\
\hline
\end{tabular}
\caption{Summary of results and provenance.}
\label{tab:summary}
\end{table}

\subsection{Setting and notation}\label{sec:setting}

The latent state is $z\in\RR^d$.
A (stochastic) world model predicts $z'=f(z,a,u)$ with driving noise $u\sim\gamma:=\Nrm(0,I_k)$ drawn independently at every step;
a deterministic model is the special case without $u$.
Fix the current latent state $s$ and a candidate set $\mathcal A$ of action sequences ($|\mathcal A|=N$; everything below holds for any $\mathcal A$ on which the stated minima are attained).
For $a\in\mathcal A$ let $Z_a=Z_H(s,a,U)$ be the terminal latent with law $P_a$ (finite second moment), mean $\mu_a$ and covariance $\Sigma_a$.
Let $g\in\RR^d$ be the goal, $c(z)=\|z-g\|^2$, let $F\subset\RR^d$ be a measurable failure set, $\kappa\ge0$, and
\[
c_\kappa(z)=\|z-g\|^2+\kappa\,\indic_F(z),\qquad p_a=\PP(Z_a\in F).
\]
For $\alpha\in[0,1)$ and an integrable random variable $Y$,
$\CVaR_\alpha(Y)=\min_{t\in\RR}\{t+(1-\alpha)^{-1}\EE(Y-t)_+\}$ \cite{rockafellar2000,rockafellar2002}; it equals $\EE[Y\mid Y\ge \VaR_\alpha(Y)]$ if $Y$ has a continuous law, equals $\EE Y$ for $\alpha=0$, equals $y$ for a constant $Y\equiv y$, and satisfies $\CVaR_\alpha(Y)\ge\EE Y$.
The planning criterion is a functional $\rho$ of the law $P_a$:
\[
\rho_{\mathrm E}(P_a)=\EE\,c_\kappa(Z_a)=:J_\kappa(a),\qquad
\rho_\alpha(P_a)=\CVaR_\alpha\bigl(c_\kappa(Z_a)\bigr).
\]
A \emph{particle planner} draws $M$ i.i.d.\ copies $Z_a^{(1)},\dots,Z_a^{(M)}$ and scores $\hat J(a)=M^{-1}\sum_j c_\kappa(Z_a^{(j)})$; a \emph{deterministic (mean-embedding) planner} uses a predictor $m(a)\in\RR^d$ and scores the single point $c_\kappa(m(a))$.
Since $\rho_{\mathrm E}(\delta_x)=\rho_\alpha(\delta_x)=c_\kappa(x)$ for a point mass, \emph{the deterministic planner does not depend on which of these criteria is intended}; this is the root of the results of Section~\ref{sec:t1}.
Ties in an $\arg\min$ are resolved arbitrarily unless stated.

%======================================================================================================
\subsection{T1: the variance gap of mean-embedding planning}\label{sec:t1}

\begin{proposition}[Bias--variance decomposition; Known]\label{prop:bv}
If $\EE\|Z\|^2<\infty$ with mean $\mu$ and covariance $\Sigma$, then $\EE\|Z-g\|^2=\|\mu-g\|^2+\tr\Sigma$.
\end{proposition}
\begin{proof}
Write $Z-g=(Z-\mu)+(\mu-g)$; the cross term has mean $2\,\EE\langle Z-\mu,\mu-g\rangle=0$ and $\EE\|Z-\mu\|^2=\tr\Sigma$.
\end{proof}

\begin{proposition}[MSE-trained predictors converge to the conditional mean; Known]\label{prop:mse}
Let $(X,Y)$ be random with $\EE\|Y\|^2<\infty$, $X=(s,a)$, $\mu(X)=\EE[Y\mid X]$ and $\Sigma(X)=\operatorname{Cov}(Y\mid X)$.
For every measurable $m$ with $\EE\|m(X)\|^2<\infty$,
\[
L(m):=\EE\|Y-m(X)\|^2=L(\mu)+\EE\|m(X)-\mu(X)\|^2,\qquad L(\mu)=\EE\,\tr\Sigma(X).
\]
Hence $\mu$ is the a.s.\ unique minimiser, and if $L(m_n)\to L(\mu)$ along a sequence of predictors then $m_n\to\mu$ in $L^2(P_X)$.
The residual loss of the best deterministic predictor equals the average noise trace.
\end{proposition}
\begin{proof}
Condition on $X$ and apply Proposition~\ref{prop:bv} with $g=m(X)$: $\EE[\|Y-m(X)\|^2\mid X]=\|\mu(X)-m(X)\|^2+\tr\Sigma(X)$; take expectations.
\end{proof}

\begin{remark}[Multi-step rollouts]\label{rem:multistep}
Proposition~\ref{prop:mse} applies to the $H$-step predictor $s,a\mapsto \EE[Z_H\mid s,a]$.
A one-step MSE predictor $\bar f(z,a)=\EE[f(z,a,U)]$ composed $H$ times equals $\EE Z_H$ only if the dynamics are affine in $z$ (Jensen's inequality otherwise), so the rollout of a deterministic model carries an additional bias.
The theorem below is stated for an \emph{arbitrary} predictor $m$ so that this bias is covered.
\end{remark}

\begin{definition}[Noise premium]\label{def:premium}
For a criterion $\rho$ and predictor $m$ put $d(a)=c_\kappa(m(a))$ (the deterministic score) and
\[
\tau_\rho(a)=\rho(P_a)-d(a)\qquad(\text{``noise premium'' of }a).
\]
For $\rho_{\mathrm E}$, $\kappa=0$ and $m=\mu$ one has $\tau_{\mathrm E}(a)=\tr\Sigma_a=:v_a$ by Proposition~\ref{prop:bv}.
Write $\osc_{\mathcal A}\tau=\max_a\tau(a)-\min_a\tau(a)$ and $\Delta_v=\osc_{\mathcal A}v$ (the \emph{variance spread}).
\end{definition}

\begin{theorem}[Regret of mean-embedding planning; New]\label{thm:regret}
Let $a_d\in\arg\min_a d(a)$ and $a^\star\in\arg\min_a\rho(P_a)$, and let $R=\rho(P_{a_d})-\rho(P_{a^\star})$ be the regret of the deterministic planner under criterion $\rho$.
\begin{enumerate}
\item[(a)] (exact) $R=[d(a_d)-d(a^\star)]+[\tau_\rho(a_d)-\tau_\rho(a^\star)]$, and the first bracket is $\le0$.
\item[(b)] (bound) $0\le R\le \tau_\rho(a_d)-\tau_\rho(a^\star)\le\osc_{\mathcal A}\tau_\rho$. In particular $R=0$ if $\tau_\rho$ is constant on $\mathcal A$ (homoscedastic noise).
\item[(c)] (tightness) For every $\Delta>0$ and $\eta\in(0,\Delta)$ there is a two-candidate instance with $\rho=\rho_{\mathrm E}$, $\kappa=0$, $m=\mu$, $\osc\tau=\Delta_v=\Delta$ and $R=\Delta-\eta$. Hence the constant $1$ in $R\le\Delta_v$ cannot be improved, and the relative regret $R/\rho(P_{a^\star})=(\Delta-\eta)/\eta$ is unbounded as $\eta\to0$.
\end{enumerate}
\end{theorem}
\begin{proof}
(a),(b): by definition $\rho(P_a)=d(a)+\tau_\rho(a)$ for every $a$, so $R=d(a_d)-d(a^\star)+\tau_\rho(a_d)-\tau_\rho(a^\star)$, and $d(a_d)\le d(a^\star)$ by the choice of $a_d$. $R\ge0$ because $a^\star$ minimises $\rho(P_a)$. The last inequality is the definition of oscillation.
(c): Take $d=1$, $g=0$. Candidate 1: $Z_1=\sqrt\Delta\,\xi$, $\xi\sim\Nrm(0,1)$ (so $\mu_1=0$, $v_1=\Delta$). Candidate 2: $Z_2\equiv\sqrt\eta$ (so $v_2=0$). With $m=\mu$ the deterministic scores are $0<\eta$, hence $a_d=1$, but $J(1)=\Delta>\eta=J(2)$, so $a^\star=2$ and $R=\Delta-\eta$, $\osc v=\Delta$.
\end{proof}

\begin{corollary}[Expected cost, with predictor error; New]\label{cor:expected}
Let $\kappa=0$, $\rho=\rho_{\mathrm E}$ and suppose $\sup_a|c(m(a))-c(\mu_a)|\le\beta$. Then $R\le\Delta_v+2\beta$; for $m=\mu$ (the infinite-data limit of MSE training, Prop.~\ref{prop:mse}), $R\le v_{a_d}-v_{a^\star}\le\Delta_v$.
\end{corollary}
\begin{proof}
$\tau_{\mathrm E}(a)=v_a+[c(\mu_a)-c(m(a))]$ and the bracket lies in $[-\beta,\beta]$, so $\osc\tau_{\mathrm E}\le\Delta_v+2\beta$; apply Theorem~\ref{thm:regret}(b).
\end{proof}

\begin{corollary}[Failure-penalised cost; New]\label{cor:fail}
Let $\rho=\rho_{\mathrm E}$, $\kappa\ge0$, $m=\mu$. Then $\tau_\kappa(a)=v_a+\kappa\,(p_a-\indic_F(\mu_a))$ and
$R\le\tau_\kappa(a_d)-\tau_\kappa(a^\star)$. If no mean lies in $F$ (the typical situation: the mean-embedding never ``sees'' the hazard), then $R\le\osc_{\mathcal A}(v+\kappa p)$.
This is tight, and the regret can grow linearly in $\kappa$: in $d=1$ with $F=[B,\infty)$, $g<B$, let $Z_1\sim\Nrm(g,s^2)$ and $Z_2\equiv g-\varepsilon$ with $\varepsilon>0$.
The deterministic planner picks candidate~1 (score $0$ versus $\varepsilon^2$) and, whenever $s^2+\kappa p_1\ge\varepsilon^2$ (so that candidate~2 is the true optimum),
\[
R(\kappa)=s^2+\kappa\,\Phibar\!\Bigl(\frac{B-g}{s}\Bigr)-\varepsilon^2 \;\xrightarrow[\kappa\to\infty]{}\;\infty .
\]
\end{corollary}
\begin{proof}
$\EE c_\kappa(Z_a)=\|\mu_a-g\|^2+v_a+\kappa p_a$ and $d(a)=\|\mu_a-g\|^2+\kappa\indic_F(\mu_a)$; subtract and use Theorem~\ref{thm:regret}.
In the example $J(1)=s^2+\kappa\,\PP(Z_1\ge B)$, $J(2)=\varepsilon^2$.
\end{proof}

\begin{corollary}[$\mathrm{CVaR}_\alpha$; New]\label{cor:cvar}
Let $\kappa=0$, $\rho=\rho_\alpha$, $m=\mu$. The deterministic planner is the same for every $\alpha$, its premium is $\tau_\alpha(a)=\CVaR_\alpha(c(Z_a))-\|\mu_a-g\|^2\ge v_a$, and $R_\alpha\le\osc_{\mathcal A}\tau_\alpha$.
In the one-dimensional Gaussian example ($Z_1\sim\Nrm(g,s^2)$, $Z_2\equiv g+\varepsilon$, $0<\varepsilon^2<s^2c_\alpha$)
\[
R_\alpha=s^2c_\alpha-\varepsilon^2,\qquad c_\alpha=1+\frac{z_\alpha\varphi(z_\alpha)}{\Phibar(z_\alpha)}\ \ge\ 1+z_\alpha^2,\qquad z_\alpha=\Phi^{-1}\!\Bigl(\tfrac{1+\alpha}{2}\Bigr),
\]
where $\varphi,\Phi$ are the standard normal density and distribution function.
Thus $c_\alpha$ increases with $\alpha$ and $c_\alpha\sim 2\ln\frac1{1-\alpha}$ as $\alpha\to1$; the CVaR-regret of the deterministic planner exceeds its expected-cost regret $s^2-\varepsilon^2$ by the factor $\approx c_\alpha$.
\end{corollary}
\begin{proof}
$\CVaR_\alpha(Y)\ge\EE Y$ gives $\tau_\alpha\ge v$. For the example, $c(Z_1)=s^2\xi^2$ with $\xi\sim\Nrm(0,1)$, so $\CVaR_\alpha(c(Z_1))=s^2\,\CVaR_\alpha(\xi^2)$.
$\xi^2$ has a continuous law with $\PP(\xi^2>z^2)=2\Phibar(z)=1-\alpha$ exactly when $z=z_\alpha$, hence $\CVaR_\alpha(\xi^2)=\frac{1}{1-\alpha}\EE[\xi^2;|\xi|>z_\alpha]=\frac{2}{1-\alpha}\int_{z_\alpha}^\infty x^2\varphi(x)\,dx=\frac{2[z_\alpha\varphi(z_\alpha)+\Phibar(z_\alpha)]}{2\Phibar(z_\alpha)}=c_\alpha$, using $\int_z^\infty x^2\varphi=z\varphi(z)+\Phibar(z)$.
The inequality $c_\alpha\ge1+z_\alpha^2$ is Mills' bound $\Phibar(z)\le\varphi(z)/z$. Candidate~2 is deterministic with cost $\varepsilon^2$ for all $\alpha$; the deterministic planner picks candidate~1 (score $0$).
\end{proof}

\begin{remark}
Theorem~\ref{thm:regret} says that the price of ignoring the predictive spread is controlled by how much the noise premium varies \emph{across the candidates being compared}.
This is why a bimodal-wind cliff (spread depends on the distance to the pit, $p_a$ varies strongly) punishes the mean-embedding planner, whereas a model with homoscedastic additive noise would not.
The regret is a statement about the objective, not about estimation: even a perfect mean predictor incurs it.
\end{remark}

%======================================================================================================
\subsection{T2: tail blindness amplified by selection}\label{sec:t2}

\paragraph{Model.}
There are $N_r\ge1$ \emph{risky} candidates and one \emph{safe} candidate.
A particle of a risky candidate costs $c_r+\kappa\xi$, with $\xi\sim\mathrm{Bernoulli}(p)$, $0<p<1$, all particles of all candidates independent (no common random numbers); the safe candidate costs $c_s>c_r$ with certainty.
Put
\[
\theta=\frac{c_s-c_r}{\kappa},\qquad k^*=\lfloor M\theta\rfloor,\qquad J_r=c_r+\kappa p,
\]
so that $J_r>c_s\iff p>\theta$ (the \emph{mistake regime}: the risky candidates are truly worse).
With $M$ particles a risky candidate has estimate $\hat J_i=c_r+\kappa K_i/M$, $K_i\sim\Bin(M,p)$ i.i.d.
Let $F:=F_{\Bin}(k^*;M,p)=\PP(K\le k^*)$.
The planner selects the minimum estimate; a tie between the safe candidate and a risky one is resolved in favour of the risky one (the opposite convention replaces $k^*$ by $\lceil M\theta\rceil-1$ everywhere).

\begin{theorem}[Selection probability; New]\label{thm:select}
\begin{enumerate}
\item[(a)] $\PP(\text{risky selected})=1-(1-F)^{N_r}$.
\item[(b)] This is nondecreasing in $N_r$ and strictly increasing when $0<F<1$ (equivalently $\theta<1$); $F\ge(1-p)^M>0$ with equality iff $k^*=0$.
\item[(c)] $1-\PP(\text{risky selected})=\exp(-\lambda N_r)$ with $\lambda=-\ln(1-F)\ge F$; the smallest $N_r$ with $\PP(\text{risky selected})\ge1-\delta$ is $\lceil\ln(1/\delta)/\lambda\rceil\approx\ln(1/\delta)/F$. In particular $\PP\to1$ geometrically as $N_r\to\infty$ for every fixed $M$.
\item[(d)] (particle budget) If $\theta<p$, then $F\le e^{-2M(p-\theta)^2}$ and so $\PP(\text{risky selected})\le N_re^{-2M(p-\theta)^2}$; thus $M\ge\ln(N_r/\delta)/(2(p-\theta)^2)$ particles suffice for error probability $\le\delta$: the required $M$ grows only logarithmically in $N_r$.
\end{enumerate}
\end{theorem}
\begin{proof}
(a) The risky candidate $i$ beats the safe one iff $c_r+\kappa K_i/M\le c_s$, i.e.\ $K_i\le M\theta$, i.e.\ $K_i\le k^*$ (integer $K_i$).
A risky candidate is selected iff at least one risky candidate does so, and the $K_i$ are independent: $\PP=1-\PP(\forall i:K_i>k^*)=1-(1-F)^{N_r}$.
(b) $x\mapsto1-(1-F)^x$ is nondecreasing and strictly increasing if $F\in(0,1)$; $F\ge\PP(K=0)=(1-p)^M$ with equality iff $k^*=0$; $F<1$ iff $k^*<M$ iff $\theta<1$.
(c) $(1-F)^{N_r}=e^{-\lambda N_r}$ and $-\ln(1-F)\ge F$; solve $e^{-\lambda N_r}\le\delta$.
(d) Since $k^*/M\le\theta<p$, Hoeffding's inequality \cite{hoeffding1963} gives $\PP(K\le k^*)=\PP(K/M-p\le-(p-k^*/M))\le e^{-2M(p-k^*/M)^2}\le e^{-2M(p-\theta)^2}$; then $1-(1-F)^{N_r}\le N_rF$.
\end{proof}

\begin{theorem}[Excess cost and optimism; New]\label{thm:excess}
In the mistake regime $p>\theta$ the best action is the safe one (true cost $c_s$), and the planner's selected plan has true cost $c_s+\kappa(p-\theta)\,\indic\{\text{risky selected}\}$. Hence
\begin{align*}
\EE[\text{true cost of selected plan}]-c_s&=\kappa(p-\theta)\bigl[1-(1-F)^{N_r}\bigr]\ \uparrow\ \kappa(p-\theta)=J_r-c_s\quad(N_r\to\infty),\\
\EE\bigl[\min_i\hat J_i\bigr]&=c_r+\frac\kappa M\sum_{k=1}^{M}\bigl(1-F_{\Bin}(k-1;M,p)\bigr)^{N_r}\ \downarrow\ c_r\quad(N_r\to\infty).
\end{align*}
So in the limit the planner is wrong with probability one, loses the full gap $J_r-c_s$ and \emph{believes} that the plan costs $c_r$, an optimism of $\kappa p$ (the winner's curse or optimiser's curse \cite{smith2006}).
\end{theorem}
\begin{proof}
All risky candidates have the same true cost $J_r$, so the true cost of the selected plan is $J_r$ if a risky plan is selected and $c_s$ otherwise; use Theorem~\ref{thm:select}(a).
For the second line, $\EE\min_iK_i=\sum_{k=1}^M\PP(\min_iK_i\ge k)=\sum_{k=1}^M(1-F_{\Bin}(k-1;M,p))^{N_r}$, each term is in $[0,1)$ and tends to $0$ because $\PP(K\le k-1)\ge\PP(K=0)>0$ for $k\ge1$.
\end{proof}

\begin{remark}
If the risky candidates have different failure probabilities $p_i$, the same proof gives $\PP(\text{risky selected})=1-\prod_i(1-F_{\Bin}(k^*;M,p_i))$.
The mechanism needs only (i) independent estimation errors and (ii) selection of an extremum; it is not specific to Bernoulli penalties.
With common random numbers across candidates the independence in (a) fails and the formula is no longer exact.
\end{remark}

\begin{theorem}[Fresh-sample verification of the top-$K$; New]\label{thm:verify}
Run stage~1 as above with $M$ particles per risky candidate and keep the set $S$ of the $K$ risky candidates with the smallest $\hat J_i$ (ties broken uniformly at random, $N_r\ge K$).
Stage~2 draws $M'$ \emph{fresh} independent particles for each $i\in S$, giving estimates $\hat J_i'$, and the final choice is the minimiser of $\{\hat J'_i:i\in S\}\cup\{c_s\}$ (ties to risky). Let $F'=F_{\Bin}(\lfloor M'\theta\rfloor;M',p)$.
\begin{enumerate}
\item[(a)] (unbiasedness) Conditionally on the stage-1 data, $\EE[\hat J'_i\mid\text{stage 1}]=J_r$ for every $i\in S$; the stage-1 estimate of a selected candidate is biased (Theorem~\ref{thm:excess}).
\item[(b)] (risk) $\PP(\text{risky selected})=1-(1-F')^K$, which does \emph{not} depend on $N_r$ or on $M$.
The excess true cost is $\kappa(p-\theta)\bigl[1-(1-F')^K\bigr]\le\kappa(p-\theta)\,K e^{-2M'(p-\theta)^2}$ (for $p>\theta$), so $M'\ge\ln(K/\delta)/(2(p-\theta)^2)$ gives risk $\le\delta$.
\item[(c)] (residual optimism) $\EE\bigl[\min_{i\in S}\hat J_i'\bigr]=c_r+\frac\kappa{M'}\sum_{k=1}^{M'}(1-F_{\Bin}(k-1;M',p))^K$; the bias is exactly $0$ for $K=1$ and is the same function as in Theorem~\ref{thm:excess} evaluated at $(K,M')$ instead of $(N_r,M)$.
\end{enumerate}
\end{theorem}
\begin{proof}
$S$ is a function of the stage-1 data and of independent tie-breaking variables, and the stage-2 particles are independent of both. Hence, conditionally on stage~1, the stage-2 failure counts of the members of $S$ are i.i.d.\ $\Bin(M',p)$, which gives (a) and the conditional law needed for (b): the final choice is risky iff $\min_{i\in S}K'_i\le\lfloor M'\theta\rfloor$, with probability $1-(1-F')^K$ whatever $S$ is. The bound follows as in Theorem~\ref{thm:select}(d); (c) is the computation of Theorem~\ref{thm:excess} with $(N_r,M)\to(K,M')$.
\end{proof}

\begin{remark}[The safe fallback must stay in the final comparison]\label{rem:fallback}
If the safe plan is only one of the $N_r+1$ candidates ranked in stage~1 and the final choice is taken among the $K$ best \emph{without} re-inserting it, verification does not protect.
Let $R_1\sim\Bin(N_r,F)$ be the number of risky candidates ranked ahead of the safe one (ties to risky) and $N_r\ge K$. Then
\[
\PP(\text{risky selected})=\PP(R_1\ge K)+\PP(R_1<K)\bigl[1-(1-F')^{K-1}\bigr]\xrightarrow[N_r\to\infty]{}1 ,
\]
because for $R_1\ge K$ the set $S$ contains only risky plans and one of them is chosen regardless of the fresh estimates. (Proof: condition on $R_1$; if $R_1<K$ the verified set contains the safe plan and exactly $K-1$ risky ones.)
Verification is therefore only useful relative to an independently evaluated alternative (an incumbent, a fallback, or an absolute threshold).
\end{remark}

\begin{proposition}[Monotone shrinkage and ranking invariance; New]\label{prop:shrink}
Let $\hat J_1,\dots,\hat J_N$ be estimates based on $M_1,\dots,M_N$ particles and let $\tilde J_i=\Phi(\hat J_i;M_i,\hat\eta)$ where $\hat\eta$ is any (data-dependent) hyper-parameter \emph{common to all candidates} and $x\mapsto\Phi(x;M,\eta)$ is strictly increasing.
\begin{enumerate}
\item[(a)] If $M_1=\dots=M_N$, then $\tilde J_i<\tilde J_j\iff\hat J_i<\hat J_j$ for all $i,j$: shrinkage cannot change the $\arg\min$, nor any ranking (so it cannot change $\PP(\text{risky selected})$ in Theorem~\ref{thm:select}).
This covers the normal--normal empirical-Bayes rule $\tilde J_i=\hat m_0+B(\hat\tau^2)(\hat J_i-\hat m_0)$, $B=\hat\tau^2/(\hat\tau^2+\sigma^2/M)\in[0,1)$ with $\hat m_0,\hat\tau^2$ estimated from the pooled data, and James--Stein-type rules \cite{efron1973} (for $B>0$; $B=0$ makes all candidates tie).
\item[(b)] If the $M_i$ differ (e.g.\ after racing), the $\arg\min$ can change. Example: $\sigma^2=4$, $\tau^2=1$, $m_0=2$, so $B(M)=\tau^2/(\tau^2+\sigma^2/M)$; candidate~1 has $\hat J_1=0.9$ from $M_1=4$, candidate~2 has $\hat J_2=1.0$ from $M_2=64$. Raw: candidate~1; shrunk: $\tilde J_1=2+\tfrac12(0.9-2)=1.45>\tilde J_2=2+\tfrac{16}{17}(1-2)=1.0588$: candidate~2.
\item[(c)] Equal $M$ is not sufficient if the shrinkage factor depends on the candidate (e.g.\ on per-candidate variance estimates): the map is then not common.
\end{enumerate}
\end{proposition}
\begin{proof}
(a) With equal $M$ every candidate is mapped by the same strictly increasing function $\Phi(\cdot;M,\hat\eta)$. (b) is arithmetic: $B(4)=1/2$, $B(64)=16/17$. (c) is exhibited numerically in Section~\ref{sec:verification}.
\end{proof}

\begin{remark}
Racing and shrinkage thus can alter the ranking only through the \emph{unequal-sample-size} channel. They do not remove the extremum-selection bias of Theorem~\ref{thm:excess}, which acts on the order statistic of the raw estimates; unbiased protection requires fresh data (Theorem~\ref{thm:verify}). Whether unequal-$M$ shrinkage helps or hurts depends on the allocation rule; we make no claim of improvement.
Pooling stage-1 and stage-2 samples into one estimate (as in a racing scheme that re-uses samples) does not restore unbiasedness, because the stage-1 part is selected on.
\end{remark}

%======================================================================================================
\subsection{T3: energy-score training, identifiability and consistency}\label{sec:t3}

For $P$ on $\RR^d$ with $\EE_P\|X\|<\infty$ and $y\in\RR^d$ the \emph{energy score} (with exponent $1$) is
$\ES(P,y)=\EE_P\|X-y\|-\tfrac12\EE_P\|X-X'\|$, $X,X'\sim P$ independent; for a point mass $\ES(\delta_m,y)=\|m-y\|$.
The \emph{energy distance} is $\mathcal E(P,Q)=2\EE\|X-Y\|-\EE\|X-X'\|-\EE\|Y-Y'\|$ ($X,X'\sim P$, $Y,Y'\sim Q$, all independent).

\begin{theorem}[Strict propriety of the energy score; Known \cite{gneiting2007,szekely2003,szekely2013}]\label{thm:propriety}
For $P,Q$ with finite first moments and $Y\sim Q$,
\[
\EE\,\ES(P,Y)-\EE\,\ES(Q,Y)=\tfrac12\,\mathcal E(P,Q),\qquad
\mathcal E(P,Q)=\frac1{c_d}\int_{\RR^d}\frac{|\phi_P(t)-\phi_Q(t)|^2}{\|t\|^{d+1}}\,dt\ \ge0,
\]
with $c_d=\pi^{(d+1)/2}/\Gamma((d+1)/2)$ and $\phi$ the characteristic function. Hence $\mathcal E(P,Q)=0\iff P=Q$, i.e.\ $\ES$ is strictly proper relative to the class of laws with finite first moment.
\end{theorem}
\begin{proof}[Proof (recalled for completeness)]
The first identity is algebra: $\EE\,\ES(Q,Y)=\EE\|Y-Y'\|-\tfrac12\EE\|Y-Y'\|=\tfrac12\EE\|Y-Y'\|$, and $\EE\,\ES(P,Y)=\EE\|X-Y\|-\tfrac12\EE\|X-X'\|$.
For the second use $\|x\|=c_d^{-1}\int(1-\cos\langle t,x\rangle)\|t\|^{-d-1}dt$ \cite{szekely2013} and $\EE\cos\langle t,X-Y\rangle=\operatorname{Re}\,\phi_P(t)\overline{\phi_Q(t)}$, $\EE\cos\langle t,X-X'\rangle=|\phi_P(t)|^2$; by Fubini (finite first moments) the integrand becomes $-2\operatorname{Re}\phi_P\bar\phi_Q+|\phi_P|^2+|\phi_Q|^2=|\phi_P-\phi_Q|^2$. A law is determined by its characteristic function.
\end{proof}

\paragraph{Noise-input models.}
Let $Q_\theta(\cdot\mid x)$ be the law of $f_\theta(x,U)$, $U\sim\gamma=\Nrm(0,I_k)$, and let data $(X,Y)\sim P_{XY}$.
Define the population risk $L(\theta)=\EE\,\ES\bigl(Q_\theta(\cdot\mid X),Y\bigr)$ and $L_\star=\tfrac12\EE\|Y-Y'\|$ ($Y'$ an independent draw from $P_{Y|X}$ given $X$).
By Theorem~\ref{thm:propriety} applied conditionally on $X$,
\begin{equation}\label{eq:risk}
L(\theta)-L_\star=\tfrac12\,\EE_X\,\mathcal E\bigl(Q_\theta(\cdot\mid X),P_{Y|X}\bigr)\ \ge0,
\end{equation}
with equality iff $Q_\theta(\cdot\mid X)=P_{Y|X}$ a.s.

\begin{proposition}[The law is identified, the noise map is not; New]\label{prop:noisemap}
(a) If $T:\RR^k\to\RR^k$ is measurable with $T_\#\gamma=\gamma$ (any orthogonal map, a sign flip, or an $x$-dependent family $T_x$), then $f_\theta(x,T_x(\cdot))$ induces the same $Q_\theta$ and the same $L$. The minimisers of $L$ are therefore never isolated points in $\theta$-space when the class is closed under such $T$: the identified object is $\theta\mapsto Q_\theta$.
(b) If $k\ge d$, every law $P$ on $\RR^d$ is $f_\#\gamma$ for some measurable $f$ (Rosenblatt transformation \cite{rosenblatt1952}; for $k<d$ a Borel isomorphism argument, cf.\ \cite{kallenberg2002}), so no structure of $f$ beyond its pushforward is determined.
(c) Everything that depends on candidates only through their marginal laws is identified: $J_\kappa(a)$, $p_a$, $\CVaR_\alpha(c_\kappa(Z_a))$ for one-step models, and, for rollouts with independent noise at each step, the law of $Z_H$ (it is determined by the one-step kernels). The \emph{coupling} induced by sharing $U$ between candidates (common random numbers) is not identified:
let $Z_1=\mu_1+s\,U$, $Z_2=\mu_2+\epsilon\,s\,U$ with $\epsilon\in\{+1,-1\}$ and $U\sim\Nrm(0,1)$; the marginal laws do not depend on $\epsilon$ but, with $m_i=\mu_i-g$,
\[
\Var\bigl(c(Z_1)-c(Z_2)\bigr)=4s^2\,(m_1-\epsilon m_2)^2 .
\]
For $m_1m_2>0$ the variance under $\epsilon=+1$ is smaller than under $\epsilon=-1$ by the factor $((m_1-m_2)/(m_1+m_2))^2$: the variance reduction promised by common random numbers is a property of the parametrisation, not of the learned law.
\end{proposition}
\begin{proof}
(a) $\EE\,\ES(Q,Y)$ depends on $f_\theta$ only through $Q_\theta$. (b) Rosenblatt: composing the inverse conditional distribution functions with the coordinates $\Phi(u_j)$ of $u$. (c) $c(Z_1)-c(Z_2)=m_1^2-m_2^2+2s\,U(m_1-\epsilon m_2)+s^2(1-\epsilon^2)U^2$ and $\epsilon^2=1$, so the variance is $4s^2(m_1-\epsilon m_2)^2$.
\end{proof}

\begin{theorem}[Consistency of the empirical energy-score minimiser; New, standard M-estimation argument \cite{jennrich1969,newey1994}]\label{thm:consistency}
Let $(X_i,Y_i)_{i\le n}$ be i.i.d.\ from $P_{XY}$ and, independently, let $U_{ij}$ ($j\le m$, $m\ge2$) be i.i.d.\ $\gamma$.
Let
\[
\hat L_n(\theta)=\frac1n\sum_{i=1}^n\Bigl[\frac1m\sum_{j}\|f_\theta(X_i,U_{ij})-Y_i\|-\frac1{2m(m-1)}\sum_{j\ne l}\|f_\theta(X_i,U_{ij})-f_\theta(X_i,U_{il})\|\Bigr].
\]
Assume (A1) $\Theta$ is a compact metric space; (A2) $\theta\mapsto f_\theta(x,u)$ is continuous for every $(x,u)$ and $(x,u)\mapsto f_\theta(x,u)$ is measurable; (A3) there is $G$ with $\|f_\theta(x,u)\|\le G(x,u)$ for all $\theta$, $\EE\,G(X,U)<\infty$ and $\EE\|Y\|<\infty$.
Let $\Theta_\star=\arg\min_\Theta L$ and let $\hat\theta_n$ be any measurable $\varepsilon_n$-minimiser of $\hat L_n$ with $\varepsilon_n\to0$.
\begin{enumerate}
\item[(a)] $L$ is continuous, $\Theta_\star\ne\emptyset$ is compact, and $\mathrm{dist}(\hat\theta_n,\Theta_\star)\to0$ a.s.
\item[(b)] (well-specified case, A4: some $\theta_0$ has $Q_{\theta_0}(\cdot\mid X)=P_{Y|X}$ a.s.) $\Theta_\star=\{\theta:Q_\theta(\cdot\mid X)=P_{Y|X}\ \text{a.s.}\}$, and the learned predictive laws converge: $\EE_X\mathcal E(Q_{\hat\theta_n}(\cdot\mid X),P_{Y|X})\to0$ a.s.
\end{enumerate}
In general $\Theta_\star$ is a set (Proposition~\ref{prop:noisemap}); $\hat\theta_n$ itself converges only when a normalisation makes the minimiser unique.
\end{theorem}
\begin{proof}
\emph{Envelope.} The second term of the summand is at most $\frac1m\sum_jG(X_i,U_{ij})$ (since $\sum_{j\ne l}(G_j+G_l)=2(m-1)\sum_jG_j$), so $|\ell_\theta|\le D:=\frac2m\sum_jG(X_i,U_{ij})+\|Y_i\|$ with $\EE D=2\EE G+\EE\|Y\|<\infty$.
\emph{Unbiasedness.} Since the $U_{ij}$ are i.i.d., $\EE[\ell_\theta\mid X_i,Y_i]=\EE\|X-y\|-\tfrac12\EE\|X-X'\|=\ES(Q_\theta(\cdot\mid X_i),Y_i)$, so $\EE\hat L_n(\theta)=L(\theta)$.
\emph{Uniform LLN.} By (A1)--(A3), dominated convergence gives continuity of $L$, and the uniform law of large numbers for continuous dominated summands \cite{jennrich1969,newey1994} gives $\sup_\theta|\hat L_n-L|\to0$ a.s.
\emph{Argmin.} $\Theta_\star$ is nonempty and compact by continuity. For $\delta>0$ let $\Theta_\delta=\{\theta:\mathrm{dist}(\theta,\Theta_\star)\ge\delta\}$ (compact); if nonempty, $\eta:=\min_{\Theta_\delta}L-\min_\Theta L>0$. Choose $n$ so large that $\sup|\hat L_n-L|<\eta/3$ and $\varepsilon_n<\eta/3$. For $\theta\in\Theta_\delta$, $\hat L_n(\theta)>\min L+2\eta/3$, whereas $\hat L_n(\hat\theta_n)\le\min\hat L_n+\varepsilon_n<\min L+2\eta/3$, so $\hat\theta_n\notin\Theta_\delta$.
(b) By~\eqref{eq:risk} and A4, $\min L=L_\star$ and $L(\theta)=L_\star$ iff $\EE_X\mathcal E=0$ iff the laws agree a.s. (as $\mathcal E\ge0$). Moreover $\EE_X\mathcal E(Q_\theta,P)=2(L(\theta)-L_\star)$ is continuous in $\theta$ and vanishes on $\Theta_\star$, so it tends to $0$ along $\hat\theta_n$.
\end{proof}

\begin{remark}[Which energy score? Practical caveat; New]\label{rem:vstat}
The summand above uses the unbiased pair average (as in our training code). The common ``V-statistic'' variant, with the double sum over all $(j,l)$ including $j=l$, multiplies the spread term by $\rho_m=(m-1)/m$ and is \emph{not} proper \cite{ferro2014}.
For the Gaussian scale family with truth $\Nrm(0,1)$ its population risk is $R_V(s)=\sqrt{2/\pi}\sqrt{1+s^2}-\rho_ms/\sqrt\pi$, minimised at $s_m^2=\rho_m^2/(2-\rho_m^2)$: $s_2=0.378$, $s_4=0.626$, $s_8=0.788$, $s_{32}=0.940$.
\end{remark}

\begin{remark}[Point-mass predictors]
Training a deterministic predictor with the energy score gives $\ES(\delta_m,y)=\|m-y\|$, i.e.\ the spatial median, whereas MSE gives the mean; for skewed noise the two differ (for $\mathrm{Exp}(1)$ noise: mean $1$, median $\ln2$). Neither is the right input for expected-cost planning unless the cost is linear.
\end{remark}

%======================================================================================================
\subsection{T4: open-loop versus closed-loop planning under process noise}\label{sec:t4}

\paragraph{Model.}
$x_{t+1}=x_t+u_t+\sigma w_t$, $t=0,\dots,H-1$, $w_t$ i.i.d.\ $\Nrm(0,1)$, $x_0$ given, goal $g$.
An \emph{open-loop} (OL) policy is a deterministic sequence $(u_t)$ (this is what a sampling planner evaluates when it rolls candidate action sequences through a noisy model); a \emph{closed-loop} (CL) policy has $u_t=\pi_t(x_0,\dots,x_t)$.
For any criterion $V_{\mathrm{OL}}(\rho)=\inf_{\mathrm{OL}}\rho(C)\ge V_{\mathrm{CL}}(\rho)=\inf_{\mathrm{CL}}\rho(C)$ since OL policies are CL policies.

\begin{proposition}[Variance growth; New]\label{prop:var}
(a) For every OL plan $\Var(x_t)=t\sigma^2$.
(b) For the feedback $u_t=\bar u_t-k(x_t-\bar x_t)$ with constant gain $k\in(0,2)$ and nominal path $\bar x$, the deviation $e_t=x_t-\bar x_t$ satisfies $e_{t+1}=(1-k)e_t+\sigma w_t$, so
$\Var(x_t)=\sigma^2\frac{1-(1-k)^{2t}}{1-(1-k)^2}\le\frac{\sigma^2}{k(2-k)}$, and for $k=1$ (deadbeat) $\Var(x_t)=\sigma^2$.
(c) For a threshold $d=B-g>0$ and the policies ``hold at $g$'': $\PP(x_t\ge B)=\Phibar\bigl(d/(\sigma\sqrt t)\bigr)\to\tfrac12$ for OL, whereas for the feedback in (b) $\PP(x_t\ge B)\le\Phibar\bigl(d\sqrt{k(2-k)}/\sigma\bigr)$ for all $t$.
\end{proposition}
\begin{proof}
(a) $x_t=\bar x_t+\sigma\sum_{s<t}w_s$. (b) $e_t=\sigma\sum_{s<t}(1-k)^{t-1-s}w_s$ is centred Gaussian with the stated variance (geometric sum). (c) Gaussian tails; the second variance is increasing in $t$ and bounded by its limit.
\end{proof}

\begin{definition}
A functional $\rho$ of nonnegative random variables is \emph{FSD-monotone} if $Y_1\le_{\mathrm{st}}Y_2\Rightarrow\rho(Y_1)\le\rho(Y_2)$, and \emph{positively homogeneous} if $\rho(\lambda Y)=\lambda\rho(Y)$ for $\lambda>0$. Examples: $\EE$, $\VaR_\alpha$, $\CVaR_\alpha$ and all spectral risk measures.
\end{definition}

\begin{theorem}[LQG: open loop costs a factor $H$; New]\label{thm:lqg}
Let $C=(x_H-g)^2$ and let $\rho$ be FSD-monotone and positively homogeneous. Then
\[
V_{\mathrm{OL}}(\rho)=H\sigma^2\,\rho(\chi^2_1),\qquad V_{\mathrm{CL}}(\rho)=\sigma^2\rho(\chi^2_1),\qquad V_{\mathrm{OL}}/V_{\mathrm{CL}}=H,
\]
attained by any OL plan with $\bar x_H=g$ and by any CL policy with $u_{H-1}=g-x_{H-1}$, respectively. For $\rho=\EE$: $V_{\mathrm{OL}}=H\sigma^2$, $V_{\mathrm{CL}}=\sigma^2$; for $\rho=\CVaR_\alpha$: $V_{\mathrm{CL}}=\sigma^2c_\alpha$ with $c_\alpha$ from Corollary~\ref{cor:cvar}. In particular, if $0<\rho(\chi_1^2)<\infty$ then for every $H\ge2$ the optimal open-loop plan is strictly worse than the optimal closed-loop policy under every such criterion, including every $\CVaR_\alpha$, and a planner that scores candidates by open-loop rollouts but is executed with replanning overestimates the achievable cost by the factor $H$.
\end{theorem}
\begin{proof}
\emph{Lower bound, CL.} For any non-anticipative policy, $x_H-g=c+\sigma w_{H-1}$ with $c=x_{H-1}+u_{H-1}-g$ measurable with respect to $\mathcal F_{H-1}$ and $w_{H-1}$ independent of $\mathcal F_{H-1}$. For fixed real $c$ and $r\ge0$,
$\PP(|c+\sigma w|\le r)=\Phi((r-c)/\sigma)-\Phi((-r-c)/\sigma)$ is maximal at $c=0$ (the integral of a symmetric unimodal density over a symmetric interval shifted by $c$ is maximal at $c=0$; \cite{anderson1955}). Hence $\PP(C\le r^2\mid\mathcal F_{H-1})\le\PP(\sigma^2w^2\le r^2)$ and, integrating, $C\ge_{\mathrm{st}}\sigma^2\chi^2_1$; FSD-monotonicity and homogeneity give $\rho(C)\ge\sigma^2\rho(\chi_1^2)$, with equality for $u_{H-1}=g-x_{H-1}$ (then $C=\sigma^2w_{H-1}^2$).
\emph{OL.} For a deterministic plan $x_H-g\sim\Nrm(m,H\sigma^2)$ with $m=\bar x_H-g$; the same argument with $c=m$ and standard deviation $\sigma\sqrt H$ gives $C\ge_{\mathrm{st}}H\sigma^2\chi^2_1$, attained at $m=0$.
\end{proof}

\begin{remark}[Entropic risk]
For $\rho_\vartheta(C)=\vartheta^{-1}\ln\EE e^{\vartheta C}$ (FSD-monotone, not homogeneous; cf.\ \cite{jacobson1973,whittle1981}) the same argument gives $V_{\mathrm{CL}}=-\frac1{2\vartheta}\ln(1-2\vartheta\sigma^2)$ and $V_{\mathrm{OL}}=-\frac1{2\vartheta}\ln(1-2\vartheta H\sigma^2)$, the latter being $+\infty$ for $\vartheta\ge1/(2H\sigma^2)$: open-loop planning loses finiteness of the risk-sensitive criterion $H$ times earlier.
\end{remark}

\paragraph{Absorbing boundary.}
Now let the walk be absorbed (stopped) at the first time $\tau\in\{1,\dots,H\}$ with $x_\tau\ge B$ (where $g<B$), controls $|u_t|\le U$, and
\[
C=(x_H-g)^2+\kappa\,\indic\{\tau\le H\},\qquad x_H=B\ \text{if absorbed}.
\]
Define the last-step function $q(y)=\EE\bigl[(\min(y+\sigma w,B)-g)^2+\kappa\,\indic\{y+\sigma w\ge B\}\bigr]$, the expected cost of landing at mean $y$ from an alive state, and the cost of being absorbed earlier, $C_k=(B-g)^2+\kappa$.

\begin{theorem}[Strict open-loop gap with an absorbing boundary; New]\label{thm:barrier}
Let $\rho=\EE$ (expected penalised cost). Then $V_{\mathrm{CL}}\le V_{\mathrm{OL}}$ with equality for $H=1$. Assume $\sigma>0$, $U>0$, that $\inf q<C_k$ (deliberately failing is not optimal; then $q$ attains its minimum $q^\star$ at some $y^\star$, because $q(y)\to\infty$ as $y\to-\infty$ and $q(y)\to C_k$ as $y\to+\infty$) and that $y^\star-U<B$ (``the best landing point is reachable from the alive region''). Then $V_{\mathrm{CL}}<V_{\mathrm{OL}}$ strictly for every $H\ge2$.
\end{theorem}
\begin{proof}
OL plans form a compact set $[-U,U]^H$ and the OL cost is continuous in the plan, so an optimal OL plan $u^\circ$ exists. Let $y=x_{H-1}$ denote the (alive part of the) state before the last step; by dynamic programming, given $x_{H-1}<B$ and control $u$ the expected cost-to-go is $q(x_{H-1}+u)$.
Define the CL policy that follows $u^\circ$ until time $H-2$ and chooses $u_{H-1}(x)\in\arg\min_{|u|\le U}q(x+u)$ (smallest minimiser, measurable) when alive. Its cost differs from that of $u^\circ$ only through
\[
\mathrm{gap}=\EE\bigl[(q(x_{H-1}+u^\circ_{H-1})-\underline q(x_{H-1}))\,\indic\{\text{alive}\}\bigr]\ge0,\qquad \underline q(x)=\min_{|u|\le U}q(x+u).
\]
For $H\ge2$ the alive part of $x_{H-1}$ has a density that is strictly positive on $(-\infty,B)$ (a killed Gaussian walk: each step has a strictly positive Gaussian kernel). If $\mathrm{gap}=0$, then $q(x+u^\circ_{H-1})=\underline q(x)$ for almost every $x<B$. Take the nonempty open interval $I=\{x<B:|x-y^\star|<U\}$ (nonempty because $y^\star-U<B$); for $x\in I$ the window $[x-U,x+U]$ contains $y^\star$, hence $\underline q(x)=q^\star$, so $q=q^\star$ on a nonempty open interval (by continuity of $q$). But $q$ is real-analytic ($q(y)=\int h(z)\,\sigma^{-1}\varphi((z-y)/\sigma)dz$ with $h$ of quadratic growth is a Gaussian convolution) and not constant ($q(y)\to\infty$ as $y\to-\infty$), contradicting the identity theorem. So $\mathrm{gap}>0$ and $V_{\mathrm{CL}}\le\mathrm{cost}(\text{modified policy})=V_{\mathrm{OL}}-\mathrm{gap}<V_{\mathrm{OL}}$.
For $H=1$ the only decision is taken at $x_0$ and OL${}=$CL. 
\end{proof}

\begin{remark}[Scope, and what is \emph{not} true]\label{rem:firstpassage}
(i) Theorem~\ref{thm:barrier} is proved for the expected penalised cost; the failure penalty $\kappa$ is itself a form of risk aversion, and Theorem~\ref{thm:lqg} covers $\CVaR$ without a barrier. For $\CVaR$ \emph{with} a barrier closed-loop optimisation is time-inconsistent and we make no claim.
(ii) Feedback bounds the \emph{spread} (Proposition~\ref{prop:var}) and the \emph{time-$t$ exceedance} probability, but \emph{not} the cumulative first-passage probability: for any stationary Gaussian closed loop the probability of ever reaching $B$ tends to $1$ as the horizon grows. In our Monte-Carlo check, holding at $g$ with the weak gain $k=0.5$ has first-passage probability below that of open loop up to $t=20$, but \emph{above} it at $t=40$ (Section~\ref{sec:verification}). Our first conjecture, that feedback keeps the cumulative failure probability bounded, was therefore false and is not claimed.
(iii) In the LQG model without barrier all candidates have the same open-loop variance $H\sigma^2$, so Theorem~\ref{thm:regret} predicts no mean-embedding regret there; the open-loop penalty is a separate effect (value of feedback), which in the barrier model interacts with the variance-dependent failure probability.
\end{remark}

%======================================================================================================
\subsection{T5: scale recalibration cannot repair a change of shape}\label{sec:t5}

Let $P$ be the model's predictive law of a scalar latent coordinate and $Q$ the true law after a distribution shift; failure is $F_B=\{z\ge B\}$ (thresholds $B$ differ across candidates, since candidates pass the hazard at different distances).
A \emph{scale recalibration} replaces $P$ by $P_s$, the law of $\mu_P+s(Z-\mu_P)$, $Z\sim P$, $s>0$ (location-scale recalibration also frees the location).

\begin{theorem}[Scale recalibration and failure events; New]\label{thm:scale}
Let $P=\Nrm(0,1)$.
\begin{enumerate}
\item[(a)] (moment blindness) If $Q$ has mean $0$ and variance $1$, every recalibration rule that uses $Q$ only through its first two moments (variance matching, Gaussian likelihood or second-moment calibration, per-dimension or full covariance matching in higher dimension) returns $s=1$, $P_s=P$, and $P(F_B)\ne Q(F_B)$ for any $B$ at which the two laws differ.
\item[(b)] (tail mismatch, one-sided) Let $Q$ be the Laplace law with variance $1$, $\Phibar_Q(B)=Q(Z\ge B)=\tfrac12e^{-\sqrt2B}$ for $B\ge0$. For every $s>0$,
$P_s(F_B)/Q(F_B)\le e^{-B^2/(2s^2)+\sqrt2B}\to0$ as $B\to\infty$.
The scale that matches $Q(F_B)$ exactly, $s_B=B/\Phibar^{-1}(\Phibar_Q(B))$, is not constant: $s_B^2\ge B/(2\sqrt2)\to\infty$, while $s_B<1$ for small $B$.
\item[(c)] (decision conflict, general form) Suppose $B\mapsto s_B$ is not constant. Then there are two thresholds $B_1,B_2>0$ and levels $\theta_1,\theta_2$ such that the decisions ``accept the candidate with hazard at $B_j$ iff its predicted failure probability is below $\theta_j$'' ($j=1,2$) are both correct under $Q$ for no scale $s$.
Concretely for $B_1=1$, $\theta_1=0.14$, $B_2=3$, $\theta_2=0.004$: the truth is accept at $B_1$ ($Q(F_{1})=0.1216<0.14$) and reject at $B_2$ ($Q(F_3)=0.00718>0.004$); the Gaussian with scale $s$ accepts at $B_1$ iff $s<s_a=0.9257$ and rejects at $B_2$ iff $s>s_b=1.1312$, and $s_b>s_a$.
\item[(d)] (error floor, two-point law) Let $Q=\frac12(\delta_{-1}+\delta_{1})$ (mean $0$, variance $1$, the shape of the benchmark's bimodal wind). For every $s>0$, $\sup_x|F_Q(x)-\Phi(x/s)|\ge\tfrac14$; in particular $\sup_B|P_s(F_B)-Q(F_B)|\ge\tfrac14$ for all scales. (For the unit-variance Laplace law the floor $\inf_s\sup_x|F_Q-\Phi(\cdot/s)|$ is $0.028$, attained near $s=0.81$.)
\end{enumerate}
The failure probability is a functional of the whole shape of the law near the hazard; scale (or any affine) recalibration fixes a two-parameter summary and leaves the shape error untouched.
\end{theorem}
\begin{proof}
(a) By construction. (b) $\Phibar(x)\le\tfrac12e^{-x^2/2}$ for $x\ge0$, so $P_s(F_B)=\Phibar(B/s)\le\tfrac12e^{-B^2/(2s^2)}$, and divide by $\tfrac12e^{-\sqrt2B}$. Applying the same bound at $s=s_B$ to $\Phibar(B/s_B)=\tfrac12e^{-\sqrt2B}$ gives $\sqrt2B\ge B^2/(2s_B^2)$, i.e.\ $s_B^2\ge B/(2\sqrt2)$. Small $B$: the entries of Table~\ref{tab:t5}.
(c) Pick $B_1,B_2$ with $s_{B_1}<s_{B_2}$ (possible because $s_B$ is not constant; the thresholds need not be ordered). $s\mapsto P_s(F_B)=\Phibar(B/s)$ is continuous and strictly increasing, so $P_s(F_{B_j})<\theta$ iff $s<s(\theta;B_j):=B_j/\Phibar^{-1}(\theta)$. Choose $\theta_1$ slightly above $Q(F_{B_1})$ (truth: accept) and $\theta_2$ slightly below $Q(F_{B_2})$ (truth: reject). Correctness at $B_1$ requires $s<s(\theta_1;B_1)$, which is slightly above $s_{B_1}$; correctness at $B_2$ requires $s>s(\theta_2;B_2)$, slightly below $s_{B_2}$. For small enough offsets $s(\theta_1;B_1)<s(\theta_2;B_2)$, so no $s$ satisfies both. The numbers are the evaluation of the two thresholds.
(d) $F_Q$ jumps from $\tfrac12$ to $1$ at $x=1$ while $x\mapsto\Phi(x/s)$ is continuous; at $x=1$ the values $F_Q(1^-)=\tfrac12$ and $F_Q(1)=1$ cannot both be within $\tfrac14$ of $\Phi(1/s)$ unless $\Phi(1/s)=\tfrac34$, in which case the distance is exactly $\tfrac14$.
The statement about failure probabilities follows since $Q(F_B)=\tfrac12\indic\{B\le1\}$ for $B>-1$ and $\sup_B$ includes $B=1$ and $B=1^+$.
\end{proof}

\begin{table}[h]
\centering\small
\begin{tabular}{rrrrr}
\hline
$B$ & $Q(Z\ge B)$ & $\Phibar(B)$ ($s=1$) & ratio & matching $s_B$\\
\hline
0.5 & 0.2465 & 0.3085 & 1.252 & 0.730\\
1.0 & 0.1216 & 0.1587 & 1.305 & 0.857\\
2.0 & 0.0296 & 0.0228 & 0.770 & 1.060\\
3.0 & 0.0072 & 0.00135 & 0.188 & 1.226\\
4.0 & 0.0018 & 0.00003 & 0.018 & 1.370\\
\hline
\end{tabular}
\caption{Gaussian model versus unit-variance Laplace truth: failure probabilities at several thresholds and the scale that would be needed at each (output of \texttt{d4\_verify.py}, T5).}
\label{tab:t5}
\end{table}

\begin{remark}[Relation to the measurements]
The empirical benchmark finds that an oracle scalar spread rescaling (true factor $1.44$) did not beat the unadapted predictor and that per-dimension rescaling did not either.
Theorem~\ref{thm:scale} shows that this is not paradoxical: the shift changes the shape of the predictive law near the hazard (bimodal wind accumulated over the horizon, interacting with the absorbing pit), a scale factor equalises at most one tail level, and the scale that is right for one candidate is wrong for another. This is an explanation consistent with, not a proof of, the observed numbers.
\end{remark}

%======================================================================================================
\subsection{Numerical verification}\label{sec:verification}

\texttt{experiments/d4\_verify.py} (single CPU core; output in the directory \texttt{results/d4/}: the printed tables in \texttt{d4\_verify.txt}, a machine-readable summary in \texttt{d4\_verify.json}) checks every statement above against closed forms, Monte Carlo and, for the barrier model, a grid dynamic program.
Table~\ref{tab:verif} lists the largest error per claim; the Monte-Carlo errors are consistent with sampling noise (two or three standard errors at most).

%%VERIF_TABLE_BEGIN
\begin{table}[h]
\centering\small
\begin{tabular}{@{}llrl@{}}
\hline
Group & Check & Max.\ error & Result\\
\hline
T1 & bias--variance identity (rel.\ MC error) & $6.1\!\times\!10^{-4}$ & pass\\
T1 & MSE fit converges to the mean (sup error, $n=10^6$) & $8.1\!\times\!10^{-3}$ & pass\\
T1 & regret identity, Thm~\ref{thm:regret}(a) & $4.4\!\times\!10^{-16}$ & pass\\
T1 & violations of $R\le\Delta_v$ (count, 12 instances) & $0$ & pass\\
T1 & MC regret vs exact (abs.) & $5.8\!\times\!10^{-3}$ & pass\\
T1 & tightness example $R=V-\eta$ (rel.) & $1.7\!\times\!10^{-3}$ & pass\\
T1 & failure-penalised regret, $\kappa\le500$ (rel.) & $7.5\!\times\!10^{-4}$ & pass\\
T1 & $c_\alpha$ closed form vs MC tail mean (rel.) & $1.6\!\times\!10^{-3}$ & pass\\
T1 & CVaR regret (rel.) & $1.6\!\times\!10^{-3}$ & pass\\
T1 & $\EE C\le\CVaR_\alpha\le\EE C/(1-\alpha)$ (violations) & $0$ & pass\\
T2 & selection probability, 40 grid cells (abs.) & $1.3\!\times\!10^{-3}$ & pass\\
T2 & monotone in $N_r$ (violations) & $0$ & pass\\
T2 & excess true cost (abs.) & $3.0\!\times\!10^{-3}$ & pass\\
T2 & $\EE\min\hat J$ formula (abs.) & $8.5\!\times\!10^{-3}$ & pass\\
T2 & Hoeffding bound (violations) & $0$ & pass\\
T2 & fresh verification, $1-(1-F')^K$ (abs.) & $1.8\!\times\!10^{-3}$ & pass\\
T2 & pooled variant, Remark~\ref{rem:fallback} (abs.) & $9.2\!\times\!10^{-4}$ & pass\\
T2 & fresh estimate unbiased given selection (abs.) & $3.6\!\times\!10^{-3}$ & pass\\
T2 & argmin changed by shrinkage, equal $M$ (fraction) & $0$ & pass\\
T2 & unequal-$M$ flip exhibited & $0$ & pass\\
T3 & ES of a Gaussian, closed form (abs.) & $5.6\!\times\!10^{-4}$ & pass\\
T3 & ES minimised at truth, $(m,s)$ grid (dist.) & $0$ & pass\\
T3 & $\EE\ES(P,Y)-\EE\ES(Q,Y)=\frac12\mathcal E$ (abs.) & $1.8\!\times\!10^{-4}$ & pass\\
T3 & point-mass ES gives median (abs.) & $6.9\!\times\!10^{-3}$ & pass\\
T3 & CRN variance $4s^2(m_1-\epsilon m_2)^2$ (rel.) & $2.6\!\times\!10^{-4}$ & pass\\
T3 & dist.\ to $\Theta_\star$ at largest $n$ (median) & $2.3\!\times\!10^{-2}$ & pass\\
T3 & V-statistic minimiser $s_m$ (abs.) & $1.0\!\times\!10^{-3}$ & pass\\
T4 & variances in Prop.~\ref{prop:var} (rel.) & $2.9\!\times\!10^{-3}$ & pass\\
T4 & values $H\sigma^2\rho$ vs $\sigma^2\rho$ (rel.) & $2.8\!\times\!10^{-3}$ & pass\\
T4 & CVaR increasing in $|m|$ (violations) & $0$ & pass\\
T4 & entropic risk, $H=1$ (rel.) & $4.8\!\times\!10^{-4}$ & pass\\
T4 & barrier DP / OL optimum vs MC (std.\ errors) & $1.9\!\times\!10^{0}$ & pass\\
T4 & min gap $V_{\rm OL}-V_{\rm CL}$, $H=2..5$ (must be $>0$) & $0.108$ & pass\\
T4 & hypotheses of Thm~\ref{thm:barrier} hold & $0$ & pass\\
T4 & exceedance probabilities, Prop.~\ref{prop:var}(c) (abs.) & $1.8\!\times\!10^{-3}$ & pass\\
T4 & FB first passage $\le$ OL for $t\le20$ (see Remark~\ref{rem:firstpassage}) & $0$ & pass\\
T5 & Laplace tails (rel.) & $1.2\!\times\!10^{-2}$ & pass\\
T5 & scales fixing both decisions (count, must be 0) & $0$ & pass\\
T5 & floor $\ge1/4$ (shortfall) & $0$ & pass\\
\hline
\end{tabular}
\caption{Largest error per checked claim (\texttt{results/d4/d4\_verify.json}). Monte-Carlo entries are sampling errors ($S=2\times10^6$ draws unless stated). Tolerances are set at about three standard errors.}
\label{tab:verif}
\end{table}
%%VERIF_TABLE_END

Three outcomes deserve mention.
\emph{(1)} The conjecture that feedback keeps the cumulative failure probability bounded was contradicted by the first-passage simulation (Remark~\ref{rem:firstpassage}) and was replaced by the exceedance statement of Proposition~\ref{prop:var}(c).
\emph{(2)} The first version of the pooled-verification check disagreed with its formula; the discrepancy was a bug in the simulator (a risky plan is selected automatically when the safe plan is not in the verified set), not in Remark~\ref{rem:fallback}, and the corrected simulation agrees.
\emph{(3)} Candidate-specific shrinkage with equal $M$ changes the argmin in about $20\%$ of random trials, which is why Proposition~\ref{prop:shrink}(a) requires a common map.
